\documentclass[12pt,compress,aspectratio=169]{beamer}
\input{../../mybeamer}

\usetikzlibrary{decorations.pathmorphing,patterns}

\title{Class 3A: Dynamics of Connected Bodies}
\subtitle{AP Physics C--Mechanics}
\input{../../me}
\input{../../mycommands}

\begin{document}

\begin{frame}
  \maketitle
\end{frame}




\begin{frame}{Multi-Body Dynamics}
  \begin{center}
    \pic{.7}{graphics/worldslongestroadtrainwithpowertrailer8}
  \end{center}
  \begin{itemize}
  \item The objects are connected by a cable or a solid linkage with negligible
    mass
  \item All objects (usually) have the same acceleration
  \item Require multiple free-body diagrams
  \end{itemize}
\end{frame}



\begin{frame}{Solving Multi-Body Problems}
  To solve a connected-bodies problem, you can follow these procedures:
  \begin{enumerate}
  \item Draw a FBD on each of the objects
  \item Sum all the forces on all the objects along the direction of motion
    \begin{itemize}
    \item Direction of motion is usually very obvious
    \item All internal forces should cancel and do not figure into the
      acceleration of the system
    \end{itemize}
  \item Compute the acceleration of the entire system using second law of motion
    \begin{itemize}
    \item Remember that (usually) every object has the same acceleration!
    \end{itemize}
  \item Go back to the FBD of each of the objects and compute the unknown
    forces (usually tension)
  \end{enumerate}
\end{frame}



%\begin{frame}{Connected Bodies: Example}
%  \textbf{Example:} A tractor-trailer pulling two trailers starts from rest
%  and accelerates with an acceleration $a$ on a straight, level road. The mass
%  of the truck (T) is \SI{5450}{\kilo\gram}, the mass of the first trailer (A)
%  is \SI{31500}{\kilo\gram}, and the mass of the second trailer (B) is
%  \SI{19600}{\kilo\gram}.
%  \begin{enumerate}
%  \item What magnitude of force must the truck generate in order to accelerate
%    the entire vehicle?
%  \item What magnitude of force must each of the trailer hitches withstand
%    while the vehicle is accelerating?
%  \end{enumerate}
%  Assume that frictional forces are negligible in comparison with the forces
%  needed to accelerate the large masses.
%\end{frame}



\begin{frame}{Different Types of Connected Bodies}
  Multiple objects pressed against one another. There may not be friction, but
  there are definitely action/reaction forces between the blocks.
  \begin{center}
    \begin{tikzpicture}[scale=.8]
      \draw[thick] (-3,0)--(4,0);
      \draw[thick] rectangle (1,1) node[midway]{$m$};
      \draw[thick] (1,0) rectangle (3,1.5) node[midway]{$M$};
      \draw[vectors] (-1.5,.5)--(0,.5) node[pos=0,left]{$\vec F_a$};
    \end{tikzpicture}
  \end{center}
  Or multiple objects stacked on top of one another. The contact surface between
  $M$ and the floor may (or may not) have friction, while the surface between
  $M$ and $m$ must have a friction coefficient $\mu$.
  \begin{center}
    \begin{tikzpicture}[scale=.8]
      \draw[thick] (-1,0)--(5.5,0);
      \draw[thick] rectangle (3,1.5) node[midway]{$M$};
      \draw[thick] (.75,1.5) rectangle (2.25,2.5) node[midway]{$m$};
      \draw[vectors] (3,.5)--(4.5,.5) node[right]{$\vec F_a$};
    \end{tikzpicture}
  \end{center}
\end{frame}



\section{Pulley Problems}

\begin{frame}{Example Problem: Atwood Machine}
  \begin{columns}
    \column{.25\textwidth}
    \centering
    \begin{tikzpicture}[scale=.7]
      \draw[ultra thick,brown] (-1,-1.6)--(-1,0);
      \draw[ultra thick,brown] (1,0)--(1,-3);
      \draw[thick,fill=gray] circle (1.05);
      \draw[thick,fill=gray!40] circle (.95);
      \draw[line width=5.5] (0,-.15)--(0,2);
      \draw[very thick] (-2,2)--(2,2);
      \fill[white] circle (.1);
      \draw[mass] (-1.5,-1.6) rectangle +(1,-1) node[midway]{$M$};
      \draw[mass] (.5,-3) rectangle +(1,-1.4) node[midway]{$m$};
    \end{tikzpicture}   

    \column{.7\textwidth}
    An \textbf{Atwood machine} is made of two objects connected by a rope that
    runs over a pulley. The pulley allows the direction of force and direction
    of motion to change between two objects.
    
    \vspace{.2in}\textbf{Example:} The object on the left has a mass of $M$ and
    the object on the right has a mass of $m$.
    \begin{enumerate}[(a)]
    \item What is the acceleration of the masses?
    \item What is the tension in the rope?
    \end{enumerate}
  \end{columns}
\end{frame}
%\begin{frame}{Example Problem: Atwood Machine}
%  An \textbf{Atwood machine} is made of two objects connected by a rope that
%  runs over a pulley. The pulley allows the direction of force and direction
%  of motion to change between two objects.
%  \begin{columns}
%    \column{.35\textwidth}
%    \centering
%    \pic1{graphics/pulley_prob_2}
%
%    \column{.65\textwidth}
%    \textbf{Example:} The object on the left has a mass of $M$ and the object
%    on the right has a mass of $m$.
%    \begin{itemize}
%    \item What is the acceleration of the masses?
%    \item What is the tension in the rope?
%    \end{itemize}
%  \end{columns}
%\end{frame}



%\begin{frame}{A Difficult Problem!}
%  \begin{columns}
%    \column{.6\textwidth}
%    \textbf{Example:} Two blocks of mass $m$ and $M$ are connected via pulley
%    with a configuration as shown on the right. The coefficient of static
%    friction between the left block and the surface is $\mu_{s,1}$, and the
%    coefficient of static friction between the right block and the surface is
%    $\mu_{s,2}$. Formulate a mathematical inequality for the condition that no
%    sliding occurs. There may be more than one inequality. 
%    
%    \column{.4\textwidth}
%    \pic{1}{graphics/pulley_prob_6}
%  \end{columns}
%\end{frame}
%
%
%
%\begin{frame}{Multiple Pulleys}
%  When there are multiple pulleys involved, we have to remember that tension
%  force is distributed evenly along the cable.
%
%  \vspace{.2in}
%  \begin{columns}
%    \column{.6\textwidth}
%    \textbf{Example:} A block of mass $m$ is pulled, via two pulleys as shown,
%    at constant velocity along a surface inclined at angle $\theta$. The
%    coefficient of kinetic friction is $\mu_k$, between block and surface.
%    Determine the pulling force $F$. Ignore the mass of the pulleys. 
%    
%    \column{.4\textwidth}
%    \pic{1}{graphics/pulley_prob_7}
%  \end{columns}  
%\end{frame}


%\begin{frame}{One More!}
%  \begin{columns}
%    \column{.75\textwidth}
%    \textbf{Example:} A block of mass $M$ is lifted at constant velocity, via
%    an arrangement of pulleys as shown. Determine the pulling force $F$. Ignore
%    the mass of the pulleys. 
%
%    \uncover<2>{
%      \vspace{.2in}\textbf{Example:} The pulling force is replaced by a $10M$
%      mass, and was let go. What are the accelerations of the $M$ and the
%      $10M$ mass?
%    }
%    
%    \column{.25\textwidth}
%    \pic{1}{graphics/pulley_prob_9}
%  \end{columns}
%\end{frame}


\begin{frame}{A More Typical Problem}
  More typically, an Atwood machine problem is one where two objects are
  sliding on a surface. These surfaces may have (or may not) have friction. In
  this example, two blocks are connected by a massless string over a
  frictionless pulley as shown in the diagram.
  \begin{center}
    \begin{tikzpicture}[scale=1.2]
      \draw[thick,brown] (-4,.4)--(.1,.4);
      \draw[thick] (0,0)--(-5.5,0) node[midway,below]{$\mu$};
      \draw[thick,fill=magenta!20] (-4,0) rectangle (-5,.75) node[midway]{$m$};
      \begin{scope}[rotate=-30,thick]
        \draw[brown] (1,.4)--(-.05,.4);
        \draw (0,0)--(3,0) node[midway,below left]{$\mu$};
        \draw[mass] (1,0) rectangle (2.5,1) node[midway]{$M$};
      \end{scope}
      \begin{scope}[rotate=-15]
        \draw[thick,fill=gray] (0,.3) circle (.15);
        \draw[thick,fill=lightgray] (0,.3) circle (.1);
        \draw[ultra thick] (0,0)--(0,.3);
        \fill (0,.3) circle (.04);
      \end{scope}
      \draw[thick,gray!70] (0,0)--(0,-1.5);
      \draw[axes] (0,-.5) arc (270:330:.5) node[midway,below]{$\phi$};
    \end{tikzpicture}
  \end{center}
  \begin{enumerate}[(a)]
  \item Determine the acceleration of the blocks.
  \item Calculate the tension in the string.
  \end{enumerate}
\end{frame}
\end{document}
